\chapter{Il territorio e le sue storie}
\label{cap:storia_territorio}

Per capire cosa si potrebbe trovare nelle Alpi venete e perché il quadro attuale è quello che è, è necessario avere una base fatta di tre strati sovrapposti che si condizionano a vicenda in modo causale: il clima, che determina l'ecologia; l'ecologia, che determina il tipo di occupazione umana possibile; e il tipo di occupazione, che determina cosa si poteva produrre, inclusa l'arte. Trattare questi livelli come descrizioni parallele indipendenti è comodo, ma nasconde le connessioni che contano. Questo capitolo li intreccia, organizzando il materiale in una sezione geografica e geologica, una climatica e una ecologica, seguite da una storia umana suddivisa in quattro periodi di peso equivalente: Paleolitico, Mesolitico, Neolitico-Eneolitico, e Bronzo-Ferro. Il peso narrativo più elevato ricade sul Tardoglaciale e sul Mesolitico, perché è in quei millenni che le oscillazioni climatiche determinavano direttamente se le Prealpi erano accessibili o no, e il tipo di occupazione umana che ne conseguiva. Il quadro d'insieme è riassunto nelle due tavole che seguono; le sezioni successive ne sviluppano il contenuto.

% ─────────────────────────────────────────────────────────────────
% Tavola 1 di 2: Pleistocene medio → Tardoglaciale
% ─────────────────────────────────────────────────────────────────
\begin{table}[p]
\small
\setlength{\extrarowheight}{2pt}
\caption{Quadro cronologico (parte prima: Pleistocene e Tardoglaciale). Date in anni prima dell'anno zero astronomico (v.\ glossario). \textit{Continua nella tavola seguente.}}
\label{tab:cronologia_a}
\begin{tabularx}{\textwidth}{>{\raggedright\arraybackslash}X
                              >{\centering\arraybackslash}p{3.2cm}
                              >{\raggedright\arraybackslash}X}
\toprule
\textbf{Geologia / Clima} & \textbf{Anni prima anno zero} & \textbf{Presenza umana in Veneto} \\
\midrule
% ── Pleistocene medio ──────────────────────────────────────────
\rowcolor{blue!8}
Pleistocene medio. Alternanza di cicli glaciali e interglaciali. Clima mediamente più freddo dell'attuale. Lessini e Berici come isole di microclima mite nella pianura. &
500.000--130.000 &
Paleolitico inferiore. \textit{Homo heidelbergensis} e forme arcaiche. Strumenti litici in selce su Lessini e Berici. Frequentazione stagionale di grotte e ripari. \\
% ── confine Pleistocene medio / Eemiano ──────────────────────
\specialrule{0.8pt}{2pt}{2pt}
% ── Eemiano ──────────────────────────────────────────────────
\rowcolor{orange!10}
Eemiano (ultimo interglaciale). Temperature di 1--2\textdegree C superiori alle attuali. Foreste temperate sull'arco alpino. &
130.000--115.000 &
Presenza umana documentata sporadicamente. Transizione verso forme neanderthaliane. \\
% ── confine Eemiano / Würm ───────────────────────────────────
\specialrule{0.8pt}{2pt}{2pt}
% ── Würm: Paleolitico medio ──────────────────────────────────
\rowcolor{cyan!8}
Würm. Avanzata progressiva dei ghiacciai alpini. Steppa-tundra in pianura; megafauna glaciale abbondante. &
115.000--38.000 &
Paleolitico medio. \textit{Homo neanderthalensis}. Riparo Tagliente, Grotta di San Bernardino, Ponte di Veia. Caccia a cervo, cavallo, orso delle caverne; mammut nelle fasi più fredde. \\
\cmidrule(lr){3-3}
% ── Würm / UMG: Paleolitico superiore ───────────────────────
\rowcolor{cyan!8}
Würm. Ultimo Massimo Glaciale (UMG, 21.000--19.000): ghiacciai fino alle pianure. Adriatico 120 m sotto l'attuale; pianura veneta estesa verso est. &
38.000--19.000 &
Paleolitico superiore. \textit{Homo sapiens}. Fumane: arte parietale in grotta (41.000 a.p.a.z.). Gruppi Epigravettiani sulla pianura e ai margini prealpini. \\
% ── confine Würm / Tardoglaciale ─────────────────────────────
\specialrule{0.8pt}{2pt}{2pt}
% ── Tardoglaciale: Dryas pre-Bølling ────────────────────────
\rowcolor{yellow!15}
Tardoglaciale iniziale. Primo ritiro dei ghiacciai, irregolare. Prealpi ancora in gran parte inaccessibili. &
19.000--12.700 &
Epigravettiano evoluto. Cacciatori sulla pianura; siti prealpini come avamposti marginali. \\
\cmidrule(lr){3-3}
% ── Bølling-Allerød ──────────────────────────────────────────
\rowcolor{yellow!15}
Bølling-Allerød: riscaldamento rapido (+5--7\textdegree C in poche centinaia di anni). Betulla e pino avanzano in quota. Prealpi accessibili. &
12.700--10.950 &
Prima colonizzazione delle quote. Riparo Dalmeri (area veneto-trentina): lastre dipinte con ocra. Val Lastaro (Asiago). Arte su supporto portatile. \\
\cmidrule(lr){3-3}
% ── Dryas Recente ────────────────────────────────────────────
\rowcolor{yellow!15}
Dryas Recente: raffreddamento brusco (10.950--9.700). I ghiacciai riavanzano parzialmente. Le quote medie diventano di nuovo ostili. &
10.950--9.700 &
I siti alpini vengono abbandonati. Il registro archeologico si svuota alle quote medie e alte. \\
\bottomrule
\end{tabularx}
\end{table}

% ─────────────────────────────────────────────────────────────────
% Tavola 2 di 2: Olocene
% ─────────────────────────────────────────────────────────────────
\begin{table}[p]
\small
\setlength{\extrarowheight}{2pt}
\caption{Quadro cronologico (parte seconda: Olocene). Date in anni prima dell'anno zero astronomico (v.\ glossario). \textit{Continua dalla tavola precedente.}}
\label{tab:cronologia_b}
\begin{tabularx}{\textwidth}{>{\raggedright\arraybackslash}X
                              >{\centering\arraybackslash}p{3.2cm}
                              >{\raggedright\arraybackslash}X}
\toprule
\textbf{Geologia / Clima} & \textbf{Anni prima anno zero} & \textbf{Presenza umana in Veneto} \\
\midrule
% ── Olocene antico: Mesolitico ───────────────────────────────
\rowcolor{green!8}
Olocene antico. Riscaldamento rapido e definitivo. Foreste avanzano verso l'alto; megafauna glaciale scompare. Lupo, lince e orso abbondanti. &
9.700--7.500 &
Mesolitico sauveterriano. Riorganizzazione forzata dall'ecologia. Riparo Battaglia (Altopiano dei Sette Comuni), Val Lastaro. Siti fino a 2.000 m. \\
\cmidrule(lr){3-3}
% ── Olocene medio: Castelnoviano / Neolitico ─────────────────
\rowcolor{green!8}
Olocene medio. Optimum climatico (7.000--6.000 a.p.a.z.): +1--2\textdegree C. Adriatico raggiunge la configurazione attuale; laguna veneziana in formazione (6.000--5.000 a.p.a.z.). &
7.500--5.250 &
Mesolitico castelnoviano; poi Neolitico (da 7.500 a.p.a.z.). Mondeval de Sora: sepoltura con ocra e conchiglie marine. Prime comunità agricole su Lessini, Berici, Euganei. \\
\cmidrule(lr){3-3}
% ── Eneolitico ───────────────────────────────────────────────
\rowcolor{green!8}
Olocene medio-recente. Clima ancora mite. Disboscamento antropico inizia a comparire nelle sequenze polliniche. &
5.250--4.250 &
Eneolitico: primi oggetti in rame. Sepolture collettive nelle grotte prealpine. Palafitte lacustri (Fimon, Lago della Costa di Arquà Petrarca). Perifericità rispetto alla Valcamonica. \\
\cmidrule(lr){3-3}
% ── Età del Bronzo ───────────────────────────────────────────
\rowcolor{green!8}
Olocene recente. Lieve raffreddamento. Pressione antropica intensa sulle foreste; espansione di prati e pascoli d'altura. &
4.250--2.750 &
Età del Bronzo. Palafitte del Garda (UNESCO). Bostel di Rotzo (Altopiano dei Sette Comuni). Prime incisioni rupestri databili in Veneto (Val d'Assa, Garda). Piroghe fluviali. \\
\cmidrule(lr){3-3}
% ── Età del Ferro ────────────────────────────────────────────
\rowcolor{green!8}
Olocene recente. Clima vicino all'attuale. &
2.750--0 &
Età del Ferro. Veneti antichi: alfabeto venetico, situle figurate, centri urbani (Padova, Este, Altino). Conquista romana (I sec.\ a.p.a.z.). \\
\bottomrule
\end{tabularx}
\end{table}

% ─────────────────────────────────────────────────────────────────
\section{Il territorio: geomorfologia e geologia di base}

Il Veneto occupa una posizione geografica di transizione tra le catene alpine e la Pianura Padana. Questa posizione si riflette in una varietà di ambienti che raramente si trova in una regione di dimensioni comparabili: le Dolomiti e le Prealpi calcaree a nord, i Colli Euganei e i Berici nel cuore della pianura, il delta del Po a est.

Dal punto di vista geologico, la distinzione più rilevante per la storia dell'arte rupestre è quella litologica. La grande maggioranza delle Prealpi venete è calcarea: calcare massiccio cretacico nelle Dolomiti, calcari marnosi e selciferi nel settore prealpino, con intercalazioni di dolomia. Queste rocce sono erodibili, solubili nelle precipitazioni acide, e producono superfici che si levigano nel tempo anziché mantenersi aspre. L'incisione è possibile, come dimostrano i siti noti, ma la durata della leggibilità è molto inferiore rispetto alle superfici vulcaniche o metamorfiche.

Fanno eccezione i Colli Euganei, un massiccio di rocce magmatiche intrusive ed effusive (trachiti, fonoliti, rioliti) di età eocenica, e alcune aree della Lessinia con affioramenti di basalto e tufo di origine vulcanica. Queste rocce dure e resistenti costituiscono superfici più favorevoli all'incisione duratura. I Colli Euganei emergono isolati dalla pianura fino a 601 metri (Monte Venda), e i Colli Berici, di natura calcarea, raggiungono i 444 metri (Monte Alonge). Entrambi i gruppi montuosi rappresentano isole geomorfologiche nel piatto paesaggio padano: nei millenni in cui la pianura era steppa glaciale o palude, queste colline erano zone di microclima più mite, con esposizioni meridionali calde, grotte nei Berici, selce di buona qualità nella Lessinia, e acqua affidabile. Hanno funzionato, per tutta la preistoria, come rifugi ecologici e concentratori di risorse, attraendo sia la fauna sia i gruppi umani che la cacciavano.

La pianura veneta, nella sua configurazione attuale, è il prodotto di millenni di sedimentazione alluvionale. Sotto gli strati storici giacciono dossi fluviali, canali abbandonati e depositi di greto che testimoniano un paesaggio molto più dinamico e acquitrinoso di quello odierno. In questo paesaggio di pianura soggetta a inondazioni stagionali, i pochissimi punti rialzati erano di estrema importanza: piccole alture naturali, creste di antico meandro, bordi di terrazza fluviale. Quegli stessi punti hanno attirato gli insediamenti preistorici per millenni, e il reticolo dei toponimi veneti conserva ancora, nella parola \textit{motta} (piccolo rilievo nel territorio piatto), la memoria di quella geografia.

La presenza di metalli e minerali utili non è irrilevante. La Lessinia e le Prealpi vicentine hanno restituito selce di ottima qualità, estratta e lavorata almeno dal Paleolitico medio. I Colli Euganei contengono rame e altri metalli che diventano rilevanti a partire dall'Eneolitico. La disponibilità di ocre e ossidi di ferro per i pigmenti rossi è attestata in diverse località prealpine.

% ─────────────────────────────────────────────────────────────────
\section{Storia climatica: dal Pleistocene all'età del Ferro}

L'arco cronologico trattato in questo capitolo copre mezzo milione di anni. In quel lasso di tempo il clima non era stabile: il Pleistocene medio e superiore sono caratterizzati da una serie di cicli glaciale-interglaciale di durata variabile tra 40.000 e 100.000 anni, scanditi da oscillazioni della temperatura media globale di 5--10\textdegree C. Nell'Italia nordorientale questo si traduceva in alternanze tra fasi di avanzata glaciale (con ghiacciai che scendevano fino alle quote pedemontane, steppa-tundra in pianura, fauna di ambiente aperto) e fasi interglaciali (con foreste che salivano in quota, temperature simili o superiori alle attuali, fauna temperata). I siti del Paleolitico inferiore e medio della Lessinia e dei Berici si distribuiscono in questo contesto oscillante: certi livelli stratigrafici, ricchi di resti di orso delle caverne e mammut, corrispondono alle fasi più fredde; altri, con cervo e rinoceronte, alle fasi meno fredde.

L'ultimo interglaciale, detto Eemiano, è collocato tra circa 130.000 e 115.000 anni prima dell'anno zero. Le temperature estive erano di 1--2\textdegree C superiori alle attuali; le foreste temperano risalivano in quota; il livello del mare era leggermente più alto di oggi. Al termine dell'Eemiano inizia la glaciazione del Würm, lunga e complessa, con fasi di intensificazione e attenuazione. Il culmine si raggiunge durante l'Ultimo Massimo Glaciale (UMG), collocato tra circa 21.000 e 19.000 anni prima dell'anno zero: in quel momento i ghiacciai alpini scendevano fino agli sbocchi delle valli in pianura; le temperature estive erano inferiori di 8--10\textdegree C rispetto a oggi; la linea degli equilibri glaciali si trovava circa 1.100 metri più in basso dell'attuale. Il livello del mare Adriatico era inferiore di circa 120 metri: non esisteva né Venezia né la sua laguna, e il Po sfociava molto più a est, verso la costa croata attuale.

Il Tardoglaciale (19.000--9.700 anni prima dell'anno zero) è il periodo di maggiore rilevanza per la preistoria veneta, e la sua struttura interna merita attenzione. Non è un semplice scioglimento graduale: è una sequenza di oscillazioni rapide. Una prima fase di riscaldamento porta al Bølling-Allerød (12.700--10.950 anni prima dell'anno zero), un periodo caldo in cui le temperature salirono di 5--7\textdegree C in poche centinaia di anni, i ghiacciai si ritirarono significativamente, e le Prealpi diventarono per la prima volta accessibili a gruppi umani. Poi, circa 10.950 anni prima dell'anno zero, il Dryas Recente: un raffreddamento brusco di analoga intensità in pochi decenni, con ghiacciai che riavanzano parzialmente e quote medie che diventano di nuovo inospitali. Il Dryas Recente dura circa 1.300 anni e si conclude circa 9.700 anni prima dell'anno zero con il passaggio all'Olocene. Il capitolo~\ref{cap:tempo_ghiacci_clima} tratta queste oscillazioni nel dettaglio; qui basta sottolineare che la finestra del Bølling-Allerød e la chiusura del Dryas Recente hanno avuto conseguenze dirette e documentabili sulla distribuzione dei siti umani nel Veneto.

Con l'inizio dell'Olocene (9.700 anni prima dell'anno zero) il riscaldamento diventa definitivo. Le temperature raggiunsero e superarono i valori attuali durante l'optimum climatico, collocato tra circa 7.000 e 6.000 anni prima dell'anno zero. Il limite superiore delle foreste si trovava alcune centinaia di metri più in alto di oggi. Il livello dell'Adriatico settentrional risalì progressivamente, raggiungendo la configurazione attuale tra 6.000 e 5.000 anni prima dell'anno zero; in quel processo si formò la laguna veneziana. I laghi prealpini, incluso il Garda, raggiunsero i livelli attuali tra la fine del Tardoglaciale e l'inizio dell'Olocene. Le torbiere dell'altopiano del Cansiglio e del Montello conservano sequenze polliniche continue che permettono di ricostruire con grande dettaglio la storia della vegetazione locale.

Nell'Olocene recente (da circa 5.000 anni prima dell'anno zero) le temperature calarono leggermente; l'età del Bronzo fu mediamente più calda dell'attuale, con il treeline più elevato; un raffreddamento progressivo tra 3.500 e 2.500 anni prima dell'anno zero contribuì all'abbandono di alcuni insediamenti di alta quota. In questo periodo si sovrappone all'effetto climatico quello della pressione antropica: le foreste erano già in ritirata per il disboscamento.

% ─────────────────────────────────────────────────────────────────
\section{Storia ecologica: paesaggi e fauna in trasformazione}

Al culmine dell'Ultimo Massimo Glaciale, il paesaggio veneto era radicalmente diverso dall'attuale. Le Prealpi e le Dolomiti erano in gran parte coperte di ghiaccio o di nuda roccia paraglaciale. La pianura sottostante era una steppa-tundra attraversata da fiumi a canali intrecciati, con scarsa copertura vegetale arborea e forte erosione eolica. La megafauna glaciale era ancora abbondante: mammut lanoso (\textit{Mammuthus primigenius}), rinoceronte lanoso (\textit{Coelodonta antiquitatis}), uro (\textit{Bos primigenius}), bisonte delle steppe (\textit{Bison priscus}), cavallo selvatico, cervo gigante (\textit{Megaloceros giganteus}) e renna. I grandi erbivori di ambiente aperto si muovevano su vasti areali stagionali nella pianura e nei fondovalle.

In questo contesto, i Colli Euganei e i Colli Berici erano anomalie significative. Emergendo di 400--600 metri dalla pianura glaciale, con esposizioni meridionali calde, grotte nei Berici, selce disponibile nella Lessinia, e fonti d'acqua affidabili, questi rilievi concentravano risorse biologiche scarse nel paesaggio circostante: vegetazione più diversificata, rifugio dal vento, microclima mite. Funzionavano come rifugi ecologici nel senso preciso del termine: aree che mantengono popolazioni di specie durante i periodi avversi per il resto del territorio. Quella stessa funzione di concentratore biologico li rendeva attraenti per i gruppi umani che cacciavano in quel paesaggio.

Con il Bølling-Allerød, il paesaggio cambiò rapidamente. La betulla e il pino avanzarono in quota, le foreste di pino si stabilirono sui versanti prealpini, e le Prealpi divennero accessibili per la prima volta. La megafauna glaciale cominciò a cedere: il mammut scomparve dall'arco alpino prima della fine del Tardoglaciale; il cervo gigante lasciò tracce sporadiche fino a circa 9.700 anni prima dell'anno zero. Il Dryas Recente provocò un parziale rollback di questi processi, con le foreste che scendevano di nuovo e la fauna glaciale che resisteva ancora in alcune aree. All'inizio dell'Olocene, con la definitiva stabilizzazione del clima, la composizione faunistica si trasformò completamente: i grandi erbivori di steppa scomparvero dal registro faunistico locale, e gli ungulati di media taglia (stambecco, camoscio, cervo, capriolo, poi alce nelle zone boschive) divennero le specie dominanti. Le foreste di quercia mista, poi di faggio, avanzarono fino alle quote alpine, spostando il treeline di 200--300 metri verso l'alto rispetto alla posizione attuale.

Un aspetto spesso trascurato è il regime di predazione. Nell'Olocene antico e medio, il lupo era abbondante sull'intero arco alpino, la lince era presente, l'orso bruno frequentava i versanti boscosi. Una popolazione di camosci o stambecchi sottoposta a predazione intensa da parte di carnivori naturali e di cacciatori umani specializzati sviluppa pattern di movimento, dimensioni del territorio e comportamenti di allarme radicalmente diversi da quelli osservabili nelle popolazioni moderne, private dei grandi carnivori per secoli. Questo ha conseguenze metodologiche per l'uso dei dati ecologici attuali nella ricostruzione del comportamento delle prede preistoriche, discusse nel capitolo~\ref{cap:cosa_cacciavano}.

Con il Neolitico e poi con il Bronzo, la deforestazione antropica aggiunse al cambiamento climatico un secondo vettore di trasformazione del paesaggio. Le sequenze polliniche mostrano aumenti marcati dei pollini di cereali e di piante ruderali a partire da circa 5.000 anni prima dell'anno zero. Il treeline scese gradualmente per effetto combinato del leggero raffreddamento e dell'espansione del pascolo d'alta quota. I prati e le brughiere che oggi caratterizzano l'Altopiano dei Sette Comuni e le zone sommitali delle Prealpi sono in parte il prodotto di questa lunga pressione antropica, non un elemento naturale del paesaggio alpino.

% ─────────────────────────────────────────────────────────────────
\section{Storia umana}

\subsection{Paleolitico (500.000--11.000 anni prima dell'anno zero)}

Le tracce più antiche di presenza umana nel Veneto risalgono al Paleolitico inferiore, tra 500.000 e 300.000 anni prima dell'anno zero, e si concentrano sui Monti Lessini e sui Colli Berici. Si tratta di strumenti in pietra scheggiata attribuibili a gruppi di \textit{Homo heidelbergensis} o forme arcaiche simili: piccoli nuclei nomadi, senza arte riconoscibile, che frequentavano grotte e ripari come basi temporanee di caccia. La concentrazione geografica su questi due rilievi non è casuale: la selce di qualità della Lessinia forniva il materiale per gli strumenti, le grotte dei Berici offrivano riparo, e il microclima più mite di entrambi i gruppi collinari era un vantaggio reale in un paesaggio di pianura glaciale. La funzione di rifugio ecologico che Euganei e Berici avranno per tutta la preistoria si afferma già in questo primo contatto documentato tra essere umano e territorio veneto.

Con il Paleolitico medio (da circa 120.000 a 38.000 anni prima dell'anno zero) la presenza umana diventa più articolata. I Neanderthal occupano Lessinia e Berici lasciando stratigrafie potenti e leggibili. I siti principali sono il Riparo Tagliente nella Valle Squaranto, la Grotta di San Bernardino nei Berici e numerosi ripari minori nella Lessinia. Tra i siti della Lessinia, Ponte di Veia merita menzione separata: si tratta di uno dei più grandi archi naturali in roccia calcarea d'Europa, nella Valle dell'Alpone, con evidenti tracce di frequentazione musteriana. Un arco naturale di quelle proporzioni non era solo un rifugio fisico: era un punto di riferimento nel paesaggio percepibile da distanze considerevoli, e probabilmente aveva un peso nella geografia mentale dei gruppi che lo frequentavano. La fauna presente nei livelli stratigrafici cambia con il clima: nei livelli corrispondenti alle fasi più fredde del Würm compaiono mammut, rinoceronte lanoso e orso delle caverne; nei livelli meno freddi predominano cervo, cavallo e camoscio. I siti della Lessinia funzionano così anche come termometri del passato, non solo come archivi di comportamento umano.

L'arrivo di \textit{Homo sapiens} in Europa è datato a circa 43.000 anni prima dell'anno zero. I Neanderthal si estinsero circa 38.000 anni prima dell'anno zero, dopo un periodo di sovrapposizione geografica discusso. La grotta di Fumane, in Lessinia (Marano di Valpolicella), è il sito dove questa transizione lascia la traccia più straordinaria del Veneto: frammenti di intonaco dipinto con figure di animali e figure ibride umano-animali, databili a circa 41.000 anni prima dell'anno zero, le testimonianze di arte parietale più antiche d'Italia. Non si tratta di arte rupestre all'aperto nel senso discusso in questo libro, ma di arte in ambiente di grotta, associata a una fase di utilizzo intensivo della caverna come base stagionale di caccia. La sua importanza sta nell'attestare che la capacità di produrre immagini simboliche era già pienamente sviluppata in quest'area alla comparsa di \textit{Homo sapiens}.

Le fasi centrali del Paleolitico superiore (Aurignaziano, Gravettiano) sono scarsamente documentate nel Veneto rispetto ad altre regioni italiane, e questa lacuna riflette probabilmente la distribuzione reale dell'occupazione: durante l'UMG (21.000--19.000 anni prima dell'anno zero), con i ghiacciai che bloccavano le valli alpine e la pianura esposta e ventosa, i cacciatori Epigravettiani vivevano principalmente sulla grande pianura padano-adriatica ora in parte sepolta sotto l'alluvio e in parte sommersa dall'Adriatico. I siti prealpini erano la periferia di quel sistema, non il suo centro. Questo spiega la scarsità di evidenze venete in quelle fasi: non c'erano meno persone, erano altrove.

Con il Bølling-Allerød (12.700--10.950 anni prima dell'anno zero) la situazione cambia in modo documentabile. Per la prima volta i ghiacciai si ritirano abbastanza da rendere le Prealpi accessibili anche alle quote medie e alte durante la stagione calda. I cacciatori seguono le prede verso l'alto. Il sito più emblematico di questa prima colonizzazione delle quote è il Riparo Dalmeri, nell'area veneto-trentina sull'Altopiano di Marcesina, datato all'Allerød e collocato a circa 1.100 metri di quota: vi sono state trovate oltre trenta lastre di calcare decorate con figure animali in ocra rossa (cervi, stambecchi, orsi, uccelli), depositate con cura nel contesto di un accampamento ricorrente. Nella stessa finestra temporale si collocano i ritrovamenti di Val Lastaro sull'Altopiano di Asiago e la sepoltura di Villabruna nel Bellunese, datata a circa 12.000 anni prima dell'anno zero. Queste non sono soste occasionali: sono campi base stagionali a cui i gruppi tornavano anno dopo anno. L'arte del Dalmeri è su lastre portatili, non su superfici fisse: un gruppo ancora semi-mobile investiva in immagini simboliche, ma su oggetti che potevano viaggiare con lui. Il passaggio all'arte su superficie fissa all'aperto richiedeva ancora qualcosa in più.

Il Dryas Recente (10.950--9.700 anni prima dell'anno zero) interrompe bruscamente questa colonizzazione. I siti di quota vengono abbandonati; il registro archeologico alle altitudini medie e alte si svuota. Non è solo un fatto climatico: è la prova che la colonizzazione delle Prealpi durante il Bølling-Allerød era contingente alle condizioni ambientali, non ancora definitiva. Il Mesolitico non sarà una continuazione lineare di quell'esperienza, ma una seconda partenza in un paesaggio trasformato.

\subsection{Mesolitico (11.700--7.500 anni prima dell'anno zero)}

La fine del Dryas Recente (9.700 anni prima dell'anno zero) segna l'inizio di una riorganizzazione ecologica che cambia profondamente le condizioni di vita dei cacciatori-raccoglitori veneti. Quella riorganizzazione non è una scelta: è una risposta forzata. Le foreste avanzano rapidamente verso le quote medie e alte, chiudendo la pianura aperta che era stata il teatro principale della vita paleolitica. Con la foresta scompaiono i grandi erbivori di ambiente aperto, cavallo e renna, che avevano costituito le prede principali per millenni. Chi rimane nel territorio veneto trova le sue prede principali nelle nuove condizioni alpine: stambecco e camoscio in quota, cervo e capriolo ai margini della foresta. Per raggiungerli sistematicamente, i cacciatori devono colonizzare le quote, e questa volta in modo definitivo.

Questa seconda colonizzazione delle Prealpi, che apre il Mesolitico, è documentata da una moltiplicazione di siti a quote progressivamente più elevate. Sull'Altopiano dei Sette Comuni (Asiago, Vicenza), il Riparo Battaglia testimonia l'occupazione sauveterriana del plateau già nelle prime fasi dell'Olocene: lo strumentario microlitico, con le sue punte a dorso e i suoi segmenti geometrici, è pensato per armi composite come la freccia, adatte alla caccia in ambienti boschivi e rupestri. Val Lastaro, sempre sull'altopiano asiaghese, ha restituito materiali sia epigravettiani sia mesolitici, confermando che queste quote erano frequentate con continuità stagionale. I cacciatori mesolitici non scoprivano per la prima volta questi luoghi: in parte li conoscevano dalle generazioni dell'Allerød. Li ritrovavano in un paesaggio trasformato e li riconfiguravano per una diversa economia.

I siti si moltiplicano alle quote di 1.500--2.000 metri e oltre. Le Prealpi vicentine e bellunesi, la Lessinia e le propaggini dell'Altopiano di Asiago sono costellate di stazioni mesolitiche di superficie, spesso individuate dai ritrovamenti di microliti geometrici caratteristici. In molti casi si tratta di posizioni con ampia visibilità sulle valli sottostanti, scelte per il controllo del territorio di caccia più che per il riparo fisico. Il modello insediativo che emerge è quello di una mobilità verticale stagionale: quote basse in inverno, quote alte in estate, con spostamenti lungo corridoi vallivi ben definiti. L'Altopiano dei Sette Comuni, con la sua posizione di cerniera tra la pianura vicentina e le valli dolomitiche, era uno di questi nodi nella rete di movimento stagionale.

Mondeval de Sora (San Vito di Cadore, Belluno), a 2.150 metri di quota nelle Dolomiti bellunesi, è il sito mesolitico veneto più noto al grande pubblico. Una sepoltura individuale con corredo ricco: utensili di osso e selce, pezzi di ornamento, ocra rossa, e conchiglie di specie marine atlantiche. Queste ultime implicano contatti a lunga distanza con le coste dell'Adriatico o del Tirreno, o scambi con gruppi in contatto con le coste. I cacciatori mesolitici dell'alta quota veneta non erano gruppi isolati: erano inseriti in reti di relazione che coprivano centinaia di chilometri. La datazione di Mondeval si colloca intorno a 8.000 anni prima dell'anno zero, nella fase castelnoviana del Mesolitico.

L'arte mesolitica all'aperto rimane la grande questione aperta del Veneto preistorico. Non esistono pannelli incisi o dipinti attribuiti con certezza al Mesolitico nella regione. Le ragioni possibili sono molteplici e non si escludono: degrado differenziale delle superfici rocciose nel corso di 9.000 anni di clima alpino; ricerca ancora insufficiente, soprattutto sulle quote medie e alte; e la possibilità, da non scartare, che i cacciatori mesolitici non praticassero l'arte rupestre su superfici fisse su larga scala, preferendo forme di espressione simbolica su supporti deperibili o portatili, come suggerisce l'esempio del Dalmeri. La questione è discussa nel capitolo~\ref{cap:tafonomia}.

\subsection{Neolitico ed Eneolitico (7.500--4.250 anni prima dell'anno zero)}

Il Neolitico veneto, datato in Italia settentrionale tra circa 7.500 e 5.250 anni prima dell'anno zero, non arriva come una sostituzione netta di popolazione. Le evidenze genetiche e culturali disponibili suggeriscono una trasformazione progressiva dei gruppi mesolitici locali, influenzati da comunità agricole che avanzavano da ovest e da est, piuttosto che una colonizzazione che cancellasse i predecessori. Il contesto climatico è favorevole: l'optimum dell'Olocene medio (7.000--6.000 anni prima dell'anno zero), con temperature leggermente superiori alle attuali e precipitazioni abbondanti, crea condizioni ottimali per l'insediamento stabile presso i bacini lacustri che in quel momento erano pieni e produttivi.

È in questa fase che il ruolo di Euganei e Berici nella storia dell'occupazione umana veneta si esprime con maggiore chiarezza. Il Lago di Fimon, nei Colli Berici a pochi chilometri da Vicenza, è uno dei siti palafitticoli neolitici meglio documentati del Veneto: un villaggio su palafitta ai margini di un bacino lacustre oggi in gran parte interrato, con ceramiche, resti botanici e faunistici che testimoniano un'economia mista di agricoltura, allevamento e caccia. Negli Euganei, il Lago della Costa presso Arquà Petrarca (Padova) ha restituito evidenze di frequentazione preistorica nelle fasi neolitica ed eneolitica: un piccolo bacino lacustre ai piedi delle colline trachitiche, sfruttato probabilmente sia per la pesca sia come punto di approdo e scambio. La continuità di questi luoghi come attrattori umani, dalle fasi paleolitiche di rifugio agli insediamenti neolitici stabili, non è casuale.

La pianura veneta neolitica non era però solo un sistema di laghi e palafitte. Era soprattutto un paesaggio di fiumi e paludi in cui i pochi punti stabili erano gli unici luoghi sicuri per insediarsi. I dossi fluviali, creste sabbiose o ghiaiose lasciate dai meandri abbandonati, erano le piattaforme naturali degli insediamenti di pianura: leggermente sopraelevati rispetto alla pianura alluvionale, con accesso diretto all'acqua, stabili in caso di piena. Queste piccole alture naturali, chiamate ancora oggi \textit{motte} nel lessico geografico veneto, erano punti di riferimento visibili nel paesaggio piatto, segnali nel territorio: non a caso molte di esse hanno restituito materiali preistorici di età diverse, a conferma di una frequentazione che attraversa i millenni. Alcune motte sono strutture artificiali, cumuli di terra costruiti intenzionalmente in epoche successive, ma poggiano spesso su rialzi naturali già usati in preistoria.

La tradizione di segnare il paesaggio con strutture permanenti si esprime nel Neolitico veneto anche attraverso alcune evidenze di strutture megalitiche, documentate in modo ancora frammentario nella fascia prealpina vicentina e sull'Altopiano dei Sette Comuni: menhir isolati, allineamenti di massi, strutture in pietra la cui attribuzione cronologica è dibattuta. La tradizione megaliitica veneta è meno sviluppata e meno studiata di quella presente in Sardegna, Puglia o in Francia, ma la sua esistenza parziale è significativa: i megaliti sono, al pari dell'arte rupestre all'aperto, una forma di marcatura permanente e visibile del paesaggio. La loro presenza periferica nel Veneto eneolitico e neolitico si inserisce nello stesso interrogativo che questo libro pone per l'arte rupestre: con quale intensità e con quale forma le comunità preistoriche venete comunicavano simbolicamente attraverso il paesaggio fisico?

L'Eneolitico (5.250--4.250 anni prima dell'anno zero) porta i primi manufatti in rame, anche se il Veneto rimane in posizione periferica rispetto ai principali circuiti di metallurgia del tempo. I confronti sono inevitabili con le regioni limitrofe: la Lombardia e il Trentino producono in questo periodo statue-stele e le incisioni della Valcamonica sono già in piena attività. In Veneto la traccia più rilevante dell'Eneolitico è funeraria: sepolture collettive nelle grotte prealpine della Lessinia e nei Berici, e sepolture in grotticella nella fascia prealpina. Il sito di Sovizzo (Vicenza) è l'unico complesso eneolitico veneto di carattere funerario-sacrale con evidenze di un certo rilievo. La perifericità del Veneto rispetto ai grandi circuiti di elaborazione culturale del tardo Eneolitico alpino è un dato di fatto: anche in questo caso, come per la successiva arte rupestre del Bronzo, la posizione geografica del Veneto orientale ai margini delle rotte di transumanza e di scambio che percorrevano il versante alpino occidentale sembra aver avuto un ruolo.

\subsection{Età del Bronzo e del Ferro (4.250--0 anni prima dell'anno zero)}

L'età del Bronzo (4.250--2.750 anni prima dell'anno zero) è il periodo meglio documentato della preistoria veneta, e il primo in cui le incisioni rupestri divengono un elemento attestato e databile con sicurezza. Il cambiamento fondamentale rispetto ai secoli precedenti è la sedentarizzazione: le comunità del Bronzo medio e recente costruiscono insediamenti stabili, usano lo stesso luogo per generazioni, e hanno interesse a segnare il territorio in modo permanente. È questa stabilità insediativa, più che un cambiamento nelle capacità tecniche o simboliche, a spiegare la comparsa dell'arte rupestre documentata nel Veneto.

I siti palafitticoli del Lago di Garda sono l'espressione più conosciuta di questa sedentarizzazione. Inclusi dal 2011 nel patrimonio mondiale dell'UNESCO come parte del sito seriale delle palafitte preistoriche circumalpine, i villaggi gardesani mostrano insediamenti stabili e tecnicamente elaborati, con bronzi lavorati, ceramiche di alta qualità e ornamenti. Sulle rocce del Garda e della Valle dell'Asse compaiono coppelle, figure geometriche e zoomorfe che trovano confronti nell'arte rupestre alpina di questa fase su entrambi i versanti della catena. La Val d'Assa, sull'Altopiano di Asiago, è la zona rupestre più rilevante del Veneto per questo periodo: un sistema di incisioni che include figure antropomorfe, scene di caccia e motivi geometrici, in un corridoio naturale che collegava l'altopiano alla pianura vicentina.

Sull'Altopiano dei Sette Comuni, il Bostel di Rotzo (Rotzo, Vicenza) è uno dei castellieri meglio conservati del Veneto: un insediamento d'altura del Bronzo finale e dell'età del Ferro, con strutture difensive in pietra, che domina visivamente una vasta porzione del territorio prealpino. I castellieri veneti, di cui il Bostel è uno degli esempi più studiati, non erano solo insediamenti: erano punti di controllo visivo del territorio, segnali di presenza permanente nel paesaggio. La loro posizione sulle creste e sugli spalloni dei pianori prealpini rispecchia una logica di organizzazione territoriale completamente diversa da quella dei cacciatori mesolitici: non mobilità verticale stagionale, ma presidio fisso di punti strategici.

La pianura veneta del Bronzo era percorsa da un sistema di vie d'acqua che integrava le rotte di terra. I fiumi erano strade: il Brenta, il Bacchiglione, l'Adige e i loro affluenti connettevano le Prealpi alla pianura e al delta del Po. Le piroghe monossili rinvenute nell'area del Bacchiglione e in altri contesti fluviali veneti attestano una navigazione fluviale organizzata, con imbarcazioni scavate da tronchi interi, alcune delle quali di dimensioni considerevoli. La mobilità del Bronzo veneto non era dunque solo terrestre: il territorio era letto e percorso anche lungo l'asse delle acque, e i siti di pianura si distribuiscono in modo coerente con questa logica, preferendo i dossi fluviali stabili vicini ai corsi d'acqua navigabili. Anche le palafitte lacustri dei bacini dei Berici ed Euganei vanno lette in questo quadro: non sono un fenomeno isolato, ma parte di una rete di insediamenti acquatici che sfruttava ogni specchio d'acqua disponibile, dai grandi laghi prealpini ai piccoli bacini collinari.

L'età del Ferro veneta (da circa 2.750 anni prima dell'anno zero) è dominata dalla cultura dei Veneti antichi, una delle grandi civiltà dell'Italia preromana. I Veneti svilupparono un sistema di scrittura (l'alfabeto venetico, derivato dall'etrusco), centri urbani complessi (Padova, Este, Altino, Vicenza), e una tradizione artistica ricca che culmina nelle situle figurate: grandi contenitori di bronzo decorati con scene narrative di banchetto, gara e processione, in uno stile narrativo senza paralleli esatti nell'Italia coeva. Le incisioni rupestri della fase venetica, dove attestate, mostrano figure di cavalieri, motivi geometrici e scene che ricorrono anche nell'arte mobiliare. La conquista romana, completata nel I secolo prima dell'anno zero, portò a una progressiva romanizzazione del paesaggio e della cultura materiale; le tradizioni religiose locali sopravvissero a lungo, come attestano i santuari di tipo \textit{fonte sacra} diffusi nell'area veneta, dove le offerte votive continuarono per secoli.

% ─────────────────────────────────────────────────────────────────
\section{Perché questo quadro conta per la domanda di questo libro}

Tre osservazioni emergono da questa rassegna e sono rilevanti per i capitoli successivi.

La prima: il Veneto preistorico non era una periferia vuota. Dalla frequentazione di \textit{Homo heidelbergensis} sui Lessini ai Veneti antichi, questo territorio è stato abitato con continuità per mezzo milione di anni. La scarsità di arte rupestre nota non rispecchia una scarsità di presenza umana: rispecchia, almeno in parte, un tipo di occupazione (mobile, stagionale, non ancora sedentaria) che non produceva marcatura permanente del paesaggio su larga scala.

La seconda: c'è una correlazione tra tipo di occupazione e tipo di produzione simbolica. I cacciatori paleolitici producono arte in grotta (Fumane) o su supporto portatile (Dalmeri). I cacciatori mesolitici probabilmente non producono arte rupestre su superficie fissa, o la producono in misura non ancora rilevata. Le prime incisioni rupestri all'aperto datate con certezza nel Veneto appartengono all'età del Bronzo, quando le comunità sono sedentarie e hanno interesse a segnare il territorio in modo permanente. Questa gradazione non è casuale: è predetta dalla logica del tipo di insediamento.

La terza: la discontinuità culturale è reale e profonda. Tra i cacciatori mesolitici del Riparo Battaglia e le comunità dell'età del Bronzo sulle rive del Garda ci sono quattro o cinque millenni di trasformazioni culturali, economiche e demografiche. Non c'è ragione di aspettarsi continuità nelle tradizioni simboliche attraverso quelle trasformazioni. Cercare arte rupestre mesolitica e arte rupestre della tarda età del Bronzo con gli stessi criteri, negli stessi luoghi, significa ignorare che erano prodotte da popolazioni con economia, mobilità e rapporto con il territorio profondamente diversi.
